\chapter{Fondamenti di computer grafica}

Per poter comprendere appieno il funzionamento della pipeline grafica, è prima necessario ripassare una serie di concetti base di algebra lineare.
Molti di questi verranno presentati in maniera semplificata, quanto basta nel nostro caso per affrontare lo studio dei prossimi argomenti.
Queste nozioni verranno poi usate per espandere le conoscenze geometriche che sono alla base della computer grafica tridimensionale.

% ====================================================
%
% 3.1 - RICHIAMI DI ALGEBRA LINEARE
%
% ====================================================
\section{Richiami di algebra lineare}
Innanzitutto si introducono velocemente le tre entità base necessarie per definire i concetti e le operazioni che verranno utilizzate.

Uno \emph{scalare} è un numero utilizzato per descrivere una grandezza.

Un \emph{punto} $P$ è un'entità la cui unica proprietà è la posizione rispetto a un sistema di riferimento.

Un \emph{vettore} $\mathbf{v}$ è un'entità caratterizzata da lunghezza (numero reale non negativo), direzione (individuata da una retta) e verso, indipendente dalla posizione nello spazio; viene rappresentato come una $n$-upla di componenti:
\[
    \mathbf{v} = \begin{bmatrix} v_1 \\ v_2 \\ v_3 \end{bmatrix}
\]

% ====================================================
% 3.1.1 - SPAZI VETTORIALI E AFFINI
% ====================================================
\subsection{Spazi vettoriali e affini}

Ora è possibile assemblare le entità base presentate per costruire delle strutture più complesse.

\begin{definizione}[Spazio vettoriale]
    Si definisce \emph{spazio vettoriale} $V$ una struttura algebrica formata da un insieme di vettori su cui sono definite l'addizione tra vettori e la moltiplicazione per uno scalare.
\end{definizione}

Poiché verranno utilizzati i vettori formati da numeri reali, introduciamo lo spazio $\mathbb{R}^n$ e un modo di rappresentare i vettori dello spazio.

\begin{definizione}[$\mathbb{R}^n$]
    Definiamo $\mathbb{R}^n$ come l'insieme dei vettori a $n$ componenti reali.
    $\mathbb{R}^n$ è uno spazio vettoriale su $\mathbb{R}$ su cui sono definite le operazioni di addizione tra due vettori $\mathbf{u}, \mathbf{v} \in \mathbb{R}^n$ e moltiplicazione di un vettore $\mathbf{v} \in \mathbb{R}^n$ per uno scalare $\lambda \in \mathbb{R}$.
    \[
        \mathbf{u} + \mathbf{v} = 
        \begin{bmatrix} u_1 + v_1 \\ \vdots \\ u_n + v_n \end{bmatrix}, 
        \qquad
        \lambda\mathbf{v} = 
        \begin{bmatrix} \lambda v_1 \\ \vdots \\ \lambda v_n \end{bmatrix}
    \]
\end{definizione}

\begin{definizione}[Combinazione lineare]
    Dati $k$ vettori $\mathbf{v}_1, \ldots, \mathbf{v}_k \in \mathbb{R}^n$ e $k$ scalari $\lambda_1, \ldots, \lambda_k \in \mathbb{R}$, si definisce \emph{combinazione lineare} la quantità:
    \[
        \sum_{i=1}^{k} \lambda_i \mathbf{v}_i = \lambda_1 \mathbf{v}_1 + \cdots + \lambda_k \mathbf{v}_k
    \]
\end{definizione}

\begin{definizione}[Base]
    Dato uno spazio vettoriale $V$, un insieme di $n$ vettori $\mathcal{B} = \{\mathbf{v}_1, \dots, \mathbf{v}_n\} \subseteq V$ è detto \emph{base} di $V$ se è linearmente indipendente e genera $V$, ovvero se qualsiasi $\mathbf{v} \in V$ è esprimibile come combinazione lineare:
    \[
        \mathbf{v} = \alpha_1 \mathbf{v}_1 + \cdots + \alpha_n \mathbf{v}_n
    \]
    dove i coefficienti $\alpha_1, \dots, \alpha_n \in \mathbb{R}$ sono unici e sono detti \emph{coordinate} di $\mathbf{v}$ rispetto alla base scelta. 
    Il numero $n$ di elementi della base è detto \emph{dimensione} di $V$.
\end{definizione}

\begin{esempio}[Base canonica di $\mathbb{R}^n$]
    La \emph{base canonica} di $\mathbb{R}^n$ è l'insieme di vettori $\mathcal{B} = \{\mathbf{e}_1, \ldots, \mathbf{e}_n\}$ dove $\mathbf{e}_i \in \mathbb{R}^n$ è il vettore con $1$ in posizione $i$-esima e $0$ altrove:
    \[
        \mathbf{e}_i = \begin{bmatrix} 0 \\ \vdots \\ 1 \\ \vdots \\ 0 \end{bmatrix} \quad \forall\, i = 1, \ldots, n
    \]
    Ad esempio, il vettore $\mathbf{v} = \begin{bmatrix}4 & 5 & 6\end{bmatrix}^T \in \mathbb{R}^3$ può essere espresso come combinazione lineare della base canonica $\mathcal{B} = \{\mathbf{e}_1, \mathbf{e}_2, \mathbf{e}_3\}$:
    \[
        \mathbf{v} = \sum_{i=1}^{3} v_i \mathbf{e}_i 
        = 4 \begin{bmatrix} 1 \\ 0 \\ 0 \end{bmatrix}
        + 5 \begin{bmatrix} 0 \\ 1 \\ 0 \end{bmatrix}
        + 6 \begin{bmatrix} 0 \\ 0 \\ 1 \end{bmatrix}
        = \begin{bmatrix} 4 \\ 5 \\ 6 \end{bmatrix}
    \]
    Le coordinate di $\mathbf{v} $ rispetto alla base canonica di $\mathbb{R}^3$ sono quindi date dal vettore $\begin{bmatrix}4 & 5 & 6\end{bmatrix}$.
\end{esempio}

Lo spazio vettoriale da solo non è sufficiente a rappresentare posizioni nello spazio: i vettori descrivono direzioni e grandezze ma non punti, rendendo necessario estendere questo concetto.

\begin{definizione}[Spazio affine]
    Uno \emph{spazio affine} è un insieme di punti che possono essere traslati da vettori di uno spazio vettoriale associato, senza che vi sia un punto di origine privilegiato.
    Lo spazio euclideo è un esempio di spazio affine reale.
\end{definizione}

Con la definizione degli spazi affini si introducono quindi nuove operazioni:
\begin{itemize}
    \item Differenza punto-punto: definisce un vettore $\mathbf{v} = P - Q$
    \item Somma punto-vettore (derivata dalla precedente): definisce un punto $P = Q + \mathbf{v}$
\end{itemize}

\begin{osservazione}
    È importante tenere a mente che:
    \begin{itemize}
        \item La somma punto-punto non è definita.
        \item I punti sono locazioni nello spazio e la differenza tra loro è data dal vettore che li congiunge.
    \end{itemize}
\end{osservazione}

% ====================================================
% 3.1.2 - COMBINAZIONI E SIMPLESSI
% ====================================================
\subsection{Combinazioni e simplessi}
Le GPU moderne sono progettate per eseguire calcoli in maniera efficiente su primitive geometriche triangolari.
In questa sezione si introducono gli strumenti matematici che generalizzano e formalizzano questi concetti.

\begin{definizione}[Combinazione affine]
    Una \emph{combinazione affine} è una combinazione lineare di punti, i cui coefficienti hanno somma $= 1$.
    \[
        P = \alpha_0 P_0 + \cdots + \alpha_n P_n = \sum_{i=0}^{n} \alpha_i P_i \quad \text{con} \quad \sum_{i=0}^{n} \alpha_i = 1
    \]
    I coefficienti $\alpha_0, \dots, \alpha_n$ sono le \emph{coordinate baricentriche} (affini) di $P$ nello spazio affine.
\end{definizione}

\begin{definizione}[Simplesso]
    Un \emph{simplesso} di dimensione $k$ è il più piccolo insieme convesso che contiene $k+1$ punti indipendenti.
    I $k+1$ punti sono detti \emph{vertici} del simplesso.
\end{definizione}

\begin{esempio}
    Un simplesso è una generalizzazione di concetti familiari come segmento, triangolo, ecc\ldots estesi a dimensioni arbitrarie.
    \begin{itemize}
        \item $k = 0$: Singolo punto
        \item $k = 1$: Segmento (due vertici distinti)
        \item $k = 2$: Triangolo (tre punti non allineati)
        \item $k = 3$: Tetraedro (quattro punti non complanari)
    \end{itemize}
\end{esempio}

\begin{definizione}[Combinazione convessa]
    Una \emph{combinazione convessa} è una combinazione affine in cui le coordinate $\in [0, 1]$:
    \[
        P = \sum_{i=0}^{n} \alpha_i P_i \quad \text{con} \quad \sum_{i=0}^{n} \alpha_i = 1\ \text{e}\ \alpha_i \in [0, 1]\ \forall\ i = 0, \ldots , n
    \]
\end{definizione}

Le coordinate baricentriche rappresentano un sistema di coordinate relative ai vertici di un simplesso.
Se vincolate a formare una combinazione convessa, garantiscono che i punti ottenibili come combinazione dei vertici di un simplesso appartengano al simplesso stesso.
Questa proprietà viene sfruttata dalle GPU per l'interpolazione degli attributi dei vertici di un triangolo, motivando la scelta del triangolo come primitiva geometrica fondamentale.

% ====================================================
% 3.1.3 - OPERAZIONI TRA VETTORI
% ====================================================
\subsection{Operazioni tra vettori} 
Si riportano ora alcune delle operazioni fondamentali effettuate con i vettori, che verranno sfruttate all'interno della pipeline grafica.

\begin{definizione}[Prodotto scalare]
    Si definisce \emph{prodotto scalare} un'applicazione $\mathbb{R}^n \times \mathbb{R}^n \to \mathbb{R}$ che gode delle seguenti proprietà $\forall\ \mathbf{u}, \mathbf{v}, \mathbf{w} \in \mathbb{R}^n$ e $\lambda \in \mathbb{R}$:
    \begin{enumerate}
        \item \emph{Simmetria}: $ \mathbf{u} \cdot \mathbf{v} = \mathbf{v} \cdot \mathbf{u}$
        \item \emph{Linearità}: $\mathbf{v} \cdot (\lambda_1\mathbf{u} + \lambda_2\mathbf{w}) = \lambda_1(\mathbf{v} \cdot \mathbf{u}) + \lambda_2(\mathbf{v} \cdot \mathbf{w})$
        \item \emph{Definitezza positiva}: $ \mathbf{v} \cdot \mathbf{v} > 0 \iff \mathbf{v} \neq 0 $
    \end{enumerate}
\end{definizione}

\begin{definizione}[Prodotto scalare standard]
    Si definisce \emph{prodotto scalare standard} $\langle \cdot, \cdot \rangle$:
    \[
        \mathbf{u}, \mathbf{v} \in \mathbb{R}^n : \langle \mathbf{u}, \mathbf{v} \rangle = \mathbf{u} \cdot \mathbf{v} = \sum_{i=1}^{n} u_i v_i = u_1 v_1 + \cdots + u_n v_n = \mathbf{u}^T \mathbf{v}
    \]
\end{definizione}

\begin{definizione}[Prodotto vettoriale]
    Dati due vettori $\mathbf{u}, \mathbf{v} \in \mathbb{R}^3$ si definisce il loro \emph{prodotto vettoriale} come il vettore $\mathbf{u} \times \mathbf{v} \in \mathbb{R}^3$ così definito:
    \begin{itemize}
        \item Modulo:
        \[
            \norm{\mathbf{u} \times \mathbf{v}} = \norm{\mathbf{u}} \cdot \norm{\mathbf{v}} \cdot \sin{\theta}
        \]
        (ove $\theta$ è l'angolo convesso fra $\mathbf{u}$ e $\mathbf{v}$)
        \item Direzione: retta perpendicolare a $\mathbf{u}$ e $\mathbf{v}$
        \item Verso: dato dalla “regola della mano destra”
    \end{itemize}
\end{definizione}

Dal momento che si lavora con vettori nello spazio, è necessario introdurre una misura per queste quantità. In computer grafica viene utilizzata la \emph{norma euclidea} (o norma 2):
\[
    \norm[2]{\mathbf{v}} = \sqrt{\langle \mathbf{v}, \mathbf{v} \rangle} = \sqrt{\mathbf{v}^T \mathbf{v}} = \sqrt{\sum_{i=1}^{n} (v_i)^2} = \sqrt{v_1^2 + \cdots + v_n^2}
\]

\begin{osservazione}
    La norma 2 è invariante rispetto alle trasformazioni ortogonali -- come rotazioni e riflessioni -- e anche rispetto alle traslazioni.
\end{osservazione}

% ====================================================
%
% 3.2 - SISTEMI DI RIFERIMENTO E TRASFORMAZIONI GEOMETRICHE
%
% ====================================================
\section{Sistemi di riferimento e trasformazioni geometriche}
Per descrivere oggetti e scene tridimensionali è necessario definire sistemi di riferimento e le trasformazioni che permettono di passare dall'uno all'altro.
In questa sezione si introduce il formalismo che consente di rappresentare uniformemente punti e vettori, e di esprimere le trasformazioni geometriche come moltiplicazioni matriciali.
In alcuni casi verrà utilizzata la notazione bidimensionale per semplicità, mantenendo invariati i concetti sottostanti.

% ====================================================
% 3.2.1 - SISTEMI DI RIFERIMENTO AFFINI E COORDINATE OMOGENEE
% ====================================================
\subsection{Sistemi di riferimento affini e coordinate omogenee}

\begin{definizione}[Sistema di riferimento]
    Un \emph{sistema di riferimento} $F$ in uno spazio affine ci consente di determinare univocamente la posizione di ogni punto dello spazio; è identificato da:
    \begin{itemize}
        \item una base di vettori linearmente indipendenti $\mathcal{B} = \{\mathbf{e}_1, \mathbf{e}_2, \dots, \mathbf{e}_n\}$;
        \item un punto di origine $P_0$.
    \end{itemize}
\end{definizione}

\begin{osservazione}
    Rappresentare punti e vettori tramite $n$ scalari (dove $n$ è la dimensione dello spazio) è ambiguo: la stessa $n$-upla potrebbe indicare sia un punto che un vettore. Estendendo convenzionalmente il prodotto scalare-punto agli spazi affini con $1 \cdot P_0 = P_0$ e $0 \cdot P_0 = 0$, possiamo scrivere:
    \[
        P = v_1 \mathbf{e}_1 + v_2 \mathbf{e}_2 + v_3 \mathbf{e}_3 + 1 \cdot P_0
    \]
    \[
        \mathbf{v} = v_1 \mathbf{e}_1 + v_2 \mathbf{e}_2 + v_3 \mathbf{e}_3 + 0 \cdot P_0
    \]
    Notiamo che la coordinata finale vale $1$ per i punti e $0$ per i vettori.
\end{osservazione}
 
\begin{definizione}[Coordinate omogenee]
    La coordinata aggiuntiva che distingue punti e vettori è detta \emph{coordinata omogenea} e indicata con $w$. Dato il sistema di riferimento $F (\mathbf{e}_1, \mathbf{e}_2, \mathbf{e}_3, P_0)$, punti e vettori in coordinate omogenee sono rappresentati da:
    \[
        P = 
            \begin{bmatrix} p_1 & p_2 & p_3 & 1 \end{bmatrix}
            \begin{bmatrix} \mathbf{e}_1 \\ \mathbf{e}_2 \\ \mathbf{e}_3 \\ P_0 \end{bmatrix}
        \qquad
        \mathbf{v} =
            \begin{bmatrix} v_1 & v_2 & v_3 & 0 \end{bmatrix}
            \begin{bmatrix} \mathbf{e}_1 \\ \mathbf{e}_2 \\ \mathbf{e}_3 \\ P_0 \end{bmatrix}
    \]
\end{definizione}

% ====================================================
% 3.2.2 - CAMBIAMENTO DI SISTEMA DI RIFERIMENTO
% ====================================================
\subsection{Cambiamento di sistema di riferimento}
\label{subsec:cambio_sdr}
Dati due sistemi di riferimento affini $F_1 = (\mathbf{x}, \mathbf{y}, O)$ e $F_2 = (\mathbf{u}, \mathbf{v}, \mathbf{e})$, rappresentiamo i vettori e l'origine di $F_2$ in funzione dei vettori e dell'origine di $F_1$:

\[
\begin{cases}
    \mathbf{u} = u_{x} \mathbf{x} + u_{y} \mathbf{y} + 0 \cdot O \\
    \mathbf{v} = v_{x} \mathbf{x} + v_{y} \mathbf{y} + 0 \cdot O \\
    \mathbf{e} = e_x \mathbf{x} + e_y \mathbf{y} + 1 \cdot O
\end{cases}
\quad \implies \quad
\begin{bmatrix}
    \mathbf{u} \\
    \mathbf{v} \\
    \mathbf{e}
\end{bmatrix} 
= 
\underbrace{
    \begin{bmatrix}
        u_x & u_y & 0 \\
        v_x & v_y & 0 \\
        e_x & e_y & 1
    \end{bmatrix}
}_{M}
\begin{bmatrix}
    \mathbf{x} \\
    \mathbf{y} \\
    O
\end{bmatrix}
\]

\begin{proposizione}
    \label{prop:cambio_sdr}
    Sia $\mathbf{a} = \begin{bmatrix} x_1 & y_1 & 1 \end{bmatrix}^T$ il vettore delle coordinate omogenee di un punto $P$ rispetto al sistema $F_1$, e sia $\mathbf{b} = \begin{bmatrix} x_2 & y_2 & 1 \end{bmatrix}^T$ il medesimo punto espresso rispetto a $F_2$. 
    Data la matrice di trasformazione delle basi $M$, valgono le seguenti relazioni di cambiamento di coordinate:
    \begin{itemize}
        \item La matrice trasposta $M^T$ mappa le coordinate dal sistema $F_2$ al sistema $F_1$:
        \[
            \mathbf{a} = M^T \mathbf{b}
        \]
        \item La matrice inversa della trasposta $(M^T)^{-1}$ mappa le coordinate dal sistema $F_1$ al sistema $F_2$:
        \[
            \mathbf{b} = (M^T)^{-1} \mathbf{a}
        \]
    \end{itemize}
\end{proposizione}

% ====================================================
% 3.2.3 - TRASFORMAZIONI GEOMETRICHE ELEMENTARI
% ====================================================
\subsection{Trasformazioni geometriche elementari}
Le trasformazioni geometriche elementari (traslazione, rotazione e scalatura) possono essere espresse in modo analogo a un cambio di sistema di riferimento.
In questo contesto, anziché cambiare gli assi di riferimento, applichiamo le matrici direttamente sulle coordinate dei punti per modificarne la posizione nello spazio.
Alla luce di quanto visto, approfondiamo come si compongono le singole matrici di trasformazione per trasformare un punto in coordinate omogenee $P = \begin{bmatrix} P_x & P_y & 1 \end{bmatrix}$ in un altro punto $P' = \begin{bmatrix} {P'}_x & {P'}_y & 1 \end{bmatrix}$, per poi visualizzarne il risultato in \cref{fig:trasformazioni_geometriche_elementari}.

\subsubsection{Traslazione}
Dato il vettore di traslazione $\mathbf{e} = \begin{bmatrix}e_x & e_y\end{bmatrix}$, otteniamo il punto $P'$ traslando $P$ del vettore $\mathbf{e}$ grazie alla matrice di traslazione $T$:
\[
P'
= 
\underbrace{
    \begin{bmatrix}
        1 & 0 & e_x \\
        0 & 1 & e_y \\
        0 & 0 & 1
    \end{bmatrix}
}_{T(\mathbf{e})}
P
\]

\subsubsection{Rotazione}
Dato l'angolo di rotazione $\theta$, otteniamo il punto $P'$ ruotando $P$ dell'angolo $\theta$ attorno all'origine del sistema di riferimento utilizzando la matrice di rotazione $R$: 
\[
P'
= 
\underbrace{
    \begin{bmatrix}
        \cos{\theta} & -\sin{\theta} & 0 \\
        \sin{\theta} & \cos{\theta} & 0 \\
        0 & 0 & 1
    \end{bmatrix}
}_{R(\theta)}
P
\]
Per le rotazioni nello spazio tridimensionale, risulta comodo utilizzare un asse di rotazione $\hat{\mathbf{n}}$ e un angolo $\theta$.

\subsubsection{Scalatura}
Dato il vettore contenente i fattori di scala $\mathbf{s} = \begin{bmatrix}s_x & s_y\end{bmatrix}$, otteniamo il punto $P'$ scalando $P$ con i fattori $\mathbf{s}$ grazie alla matrice di scala $S$:
\[
P'
= 
\underbrace{
    \begin{bmatrix}
        s_x & 0 & 0 \\
        0 & s_y & 0 \\
        0 & 0 & 1
    \end{bmatrix}
}_{S(\mathbf{s})}
P
\]
Tenendo conto che:
\begin{itemize}
    \item Fattori $s < 1$: i punti si avvicinano all'origine
    \item Fattori $s > 1$: i punti si allontanano dall'origine
    \item Se $s_x \neq s_y$: le proporzioni non sono mantenute (\emph{scalatura non uniforme})
    \item Se $s_x = s_y$: le proporzioni sono mantenute (\emph{scalatura uniforme})
\end{itemize}

\begin{figure}[H]
    \centering
    \begin{subfigure}[b]{0.32\textwidth}
        \centering
        \includegraphics[width=\textwidth]{cap3/Traslazione}
        \caption{Traslazione}
    \end{subfigure}
    \hfill
    \begin{subfigure}[b]{0.32\textwidth}
        \centering
        \includegraphics[width=\textwidth]{cap3/Rotazione}
        \caption{Rotazione}
    \end{subfigure}
    \hfill
    \begin{subfigure}[b]{0.32\textwidth}
        \centering
        \includegraphics[width=\textwidth]{cap3/Scala}
        \caption{Scalatura}
    \end{subfigure}
    \caption[Applicazione trasformazioni geometriche elementari]{Applicazione delle trasformazioni geometriche elementari a dei punti del piano}
    \label{fig:trasformazioni_geometriche_elementari}
\end{figure}

% ====================================================
%
% 3.3 - LA PIPELINE GRAFICA
%
% ====================================================
\section{La pipeline grafica}
La pipeline grafica è la sequenza di operazioni atte a costruire un'immagine partendo dagli oggetti tridimensionali presenti nella scena.
L'immagine prodotta è di tipo \emph{raster}, poiché è composta da una matrice di punti -- detti picture elements (pixel) -- ognuno contenente le proprie informazioni sul colore.

La GPU gestisce il processo di rendering tramite alcuni buffer, che non sono altro che matrici in cui ad ogni pixel è associata una cella di memoria.
I principali sono il \emph{color buffer}, che contiene il colore di ogni pixel dell'immagine finale, e il \emph{depth buffer} (o \emph{z-buffer}), che memorizza la profondità del frammento più vicino all'osservatore per ogni pixel, permettendo di gestire correttamente la visibilità tra oggetti sovrapposti.
L'insieme di questi buffer forma il \emph{framebuffer}, che viene letto e scritto nelle varie fasi della pipeline.

L'intero processo di visualizzazione può essere riassunto in quattro sottosistemi:
\begin{itemize}
    \item Applicazione: l'applicazione con cui interagisce l'utente produce le primitive geometriche tridimensionali.
    \item Geometria: le primitive vengono sottoposte a una sequenza di trasformazioni, passando da un sistema di riferimento a un altro premoltiplicando per apposite matrici, che permettono di passare dai vertici della scena 3D a coordinate normalizzate.
    \item Rasterizzazione: converte le primitive geometriche risultanti in  frammenti da visualizzare sullo schermo tramite appositi algoritmi.
    \item Visualizzazione: il dispositivo di output (ad esempio un monitor) mostra l'immagine creata.
\end{itemize}

Durante le fasi di geometria e rasterizzazione, a partire da OpenGL 2.0, il programmatore può intervenire all'interno della pipeline tramite dei programmi scritti in linguaggio OpenGL Shading Language (GLSL), che prendono il nome di \emph{shader}.
Esistono vari tipi di shader che agiscono in momenti distinti, ma gli obbligatori sono il \emph{vertex shader}, eseguito per ogni vertice all'inizio del sottosistema geometrico, e il \emph{fragment shader}, eseguito per ogni frammento prodotto dalla rasterizzazione.
Nel progetto viene utilizzato anche il \emph{geometry shader}, che viene eseguito subito dopo il vertex shader e, a differenza di quest'ultimo, riceve in input una primitiva geometrica completa, potendo così accedere a tutti i suoi vertici contemporaneamente.
In \cref{fig:pipeline} è possibile visualizzare dove intervengono durante il processo di produzione delle immagini.

\begin{figure}[ht]
    \centering
    \begin{tikzpicture}[
        box/.style={rectangle, draw, minimum width=3cm, minimum height=1.7cm, align=center},
        arrow/.style={->, thick, >={Stealth[length=8pt, width=6pt]}}
    ]
        % Nodi
        \node[box] (app) {Applicazione};
        \node[box, right=4.5cm of app] (geo) {Geometria \\ {\small (Vertex e Geometry shader)}};
        \node[box, below=1.5cm of geo] (rast) {Rasterizzazione \\ {\small (Fragment shader)}};
        \node[box, below=1.5cm of app] (visual) {Visualizzazione};

        % Frecce
        \draw[arrow] (app) -- node[above] {Vertici e primitive} (geo);
        \draw[arrow] (geo) -- node[left] {Vertici e primitive trasformati} (rast);
        \draw[arrow] (rast) -- node[above] {Pixel} (visual);
    \end{tikzpicture}
    \caption{Schema del processo di produzione di un'immagine}
    \label{fig:pipeline}
\end{figure}

Poiché nelle sezioni precedenti abbiamo studiato e analizzato gli strumenti per manipolare punti e vettori, vediamo come essi vengono sfruttati all'interno della pipeline grafica: per il sottosistema geometrico analizzeremo ogni trasformazione di sistema di riferimento, mentre per quello di rasterizzazione descriveremo il processo di conversione delle primitive in frammenti.

% ====================================================
% 3.3.1 - MODELLAZIONE
% ====================================================
\subsection{Modellazione}
Inizialmente, ogni oggetto è modellato in un proprio sistema di riferimento, detto \emph{Object Coordinate System} (OCS): vengono poi posizionati nella scena tridimensionale, il cui sistema di riferimento è detto \emph{World Coordinate System} (WCS).
Questa fase è detta di modellazione, e avviene applicando a ogni vertice le trasformazioni studiate, ovvero scalatura, rotazione e infine traslazione.
Poiché le trasformazioni geometriche sono associative ma non commutative, è necessario che vengano applicate nell'ordine corretto.
Dal momento che nelle moltiplicazioni matriciali l'ultima matrice è la prima ad agire, la matrice di modellazione è composta nella seguente maniera:
\[
    M = T(e) \cdot R(\theta) \cdot S(s)
\]

% ====================================================
% 3.3.2 - VISTA
% ====================================================
\subsection{Vista}
In questa fase, viene posizionata e orientata nella scena una telecamera virtuale.
La trasformazione di vista è una matrice che mappa tutti i vertici dal WCS al \emph{View Coordinate System} (VCS), il cui punto di origine è la posizione della telecamera (che convenzionalmente punta verso $-z$).

Uno dei modelli di camera virtuale più intuitivi e utilizzati, è quello della camera \emph{look at}, che si basa su quattro informazioni:
\begin{itemize}
    \item Punto C: posizione della telecamera nel WCS;
    \item Punto A: punto verso cui la telecamera guarda;
    \item \emph{Field of View} (FOV): angolo di visibilità verticale nella scena, ovvero quanta porzione di scena viene inquadrata;
    \item \emph{Depth of Field}: determina la distanza tra l'oggetto più vicino e quello più lontano che appaiono a fuoco.
\end{itemize}
Da queste informazioni possiamo ricavare i versori che definiscono il sistema di riferimento VCS.

La direzione di vista è il versore $\hat{\mathbf{w}}$ che da C punta verso A:
\[
    \hat{\mathbf{w}} = \frac{C - A}{\norm{C - A}}
\]

Il \emph{View Up Vector} ($VUP$) è il vettore che punta verso l'alto della telecamera.

L'asse $\hat{\mathbf{u}}$ è il versore che punta alla destra dell'osservatore: è perpendicolare al piano individuato da $VUP$ e $\hat{\mathbf{w}}$:
\[
    \hat{\mathbf{u}} = \frac{VUP \times \hat{\mathbf{w}}}{\norm{VUP \times \hat{\mathbf{w}}}}
\]

L'asse $\hat{\mathbf{v}}$ è a sua volta perpendicolare a $\hat{\mathbf{w}}$ e $\hat{\mathbf{u}}$, ma dato che sono già versori non è necessario rieseguire la normalizzazione:
\[
    \hat{\mathbf{v}} = \hat{\mathbf{w}} \times \hat{\mathbf{u}}
\]

Definiamo quindi il sistema di riferimento in coordinate di vista $VCS (\hat{\mathbf{u}}, \hat{\mathbf{v}}, \hat{\mathbf{w}}, C)$, ovvero i tre assi appena trovati con la telecamera all'origine. 
Conoscendo questi vettori risulta semplice calcolare la matrice di vista.
Come visto nella \cref{subsec:cambio_sdr} per il cambio di sistema di riferimento, la matrice di trasformazione delle basi che esprime il $VCS (\hat{\mathbf{u}}, \hat{\mathbf{v}}, \hat{\mathbf{w}}, C)$ in termini del $WCS (\hat{\mathbf{x}}, \hat{\mathbf{y}}, \hat{\mathbf{z}}, O)$ è:
\[
M =
\begin{bmatrix}
    u_x & u_y & u_z & 0 \\
    v_x & v_y & v_z & 0 \\
    w_x & w_y & w_z & 0 \\
    C_x & C_y & C_z & 1
\end{bmatrix}
\]

Poiché la nostra matrice di vista $V$ dovrà trasformare i punti da WCS a VCS, per la \cref{prop:cambio_sdr} è composta come:
\[
V =
(M^T)^{-1}
\]

% ====================================================
% 3.3.3 - PROIEZIONE
% ====================================================
\subsection{Proiezione}
In generale, la proiezione è una trasformazione geometrica che ha come dominio uno spazio di dimensione $n$ e come codominio uno spazio di dimensione $n-1$.
Una proiezione di un oggetto 3D è definita da un insieme di rette di proiezione dette proiettori che, partendo da un centro di proiezione comune, passano per tutti i punti dell'oggetto e intersecano un piano di proiezione per formare l'immagine bidimensionale.
Quando i proiettori sono rette e la proiezione avviene su un piano si parla di proiezioni geometriche piane, che vengono divise in due tipologie: le proiezioni prospettiche e le proiezioni parallele.

L'operazione di proiezione nella pipeline grafica viene divisa in due fasi: la determinazione del volume di vista -- ovvero la regione di spazio compresa tra due piani, detti piano vicino e piano lontano, che la camera virtuale è in grado di inquadrare in un dato momento -- e la sua mappatura in un cubo canonico (trattata in questa sezione), e la riduzione da spazio 3D a 2D (trattata nella sezione successiva).

In questa fase, grazie alla matrice di proiezione $P$, si passa dal VCS al sistema \emph{Normalized Device Coordinates} (NDC), un cubo canonico dai lati di lunghezza 2 centrato nell'origine tale che $x, y, z \in [-1, 1]$, così da semplificare sensibilmente i calcoli svolti negli stadi successivi.
Al termine della moltiplicazione per la matrice di proiezione, la  pipeline riporta automaticamente i punti in coordinate cartesiane (quindi con la quarta coordinata $w = 1$)  dividendo tutte le componenti per $w$, operazione detta \emph{perspective divide}.

\subsubsection{Proiezioni parallele}
Il nome proiezioni parallele deriva dalla proprietà di parallelismo dei proiettori: non c'è un centro di proiezione, ma si parla di direzione di proiezione. 
Poiché vengono mantenute le distanze, linee parallele nella realtà lo sono anche nel volume proiettato.

Si dividono a loro volta in due classi, le proiezioni oblique e quelle ortografiche.
Nelle proiezioni oblique i proiettori possono formare un angolo qualsiasi col piano di proiezione: ne è un esempio la cavaliera, utilizzata nel disegno tecnico, in cui la direzione di proiezione forma un angolo di 45 gradi con il piano di proiezione.
Nelle proiezioni ortografiche invece, la direzione di proiezione coincide con la normale al piano di proiezione: in questo caso il volume di vista corrisponde con un parallelepipedo (visibile in \cref{fig:proiezione_ortografica}), che risulta semplice da trasformare nel cubo canonico poiché è sufficiente traslarne il centroide nell'origine per poi scalarne i vertici in $[-1, 1]$.
Dati i vertici del parallelepipedo, il suo centroide è identificato dal punto:
\[
    C = 
    \begin{bmatrix} C_x & C_y & C_z \end{bmatrix} =
    \begin{bmatrix}
        \frac{left + right}{2} &
        \frac{top + bottom}{2} &
        \frac{-near - far}{2}
    \end{bmatrix}
\]
Per traslarlo nell'origine, il vettore di traslazione è semplicemente:
\[
    \mathbf{e} = O - C = \begin{bmatrix} -C_x & -C_y & -C_z \end{bmatrix}
\]
Mentre i fattori di scala per ottenere lati di lunghezza 2 saranno:
\[
    \mathbf{s} = 
    \begin{bmatrix}
        \frac{2}{right - left} &
        \frac{2}{top - bottom} &
        \frac{2}{-near + far}
    \end{bmatrix}
    = \begin{bmatrix} s_x & s_y & s_z \end{bmatrix}
\]
Otteniamo così la matrice di proiezione ortografica come la composizione della matrice di scala e di traslazione:
\[
    P_{ortho} =
    \underbrace{
        \begin{bmatrix}
            s_x &  &  &  \\
             & s_y &  &  \\
             &  & s_z &  \\
             &  &  & 1
        \end{bmatrix}
    }_{S(\mathbf{s})}
    \cdot
    \underbrace{
        \begin{bmatrix}
            1 &  &  & -C_x \\
             & 1 &  & -C_y \\
             &  & 1 & -C_z \\
             &  &  & 1
        \end{bmatrix}
    }_{T(\mathbf{e})}
    =
    \begin{bmatrix}
        s_x &  &  & -C_x s_x \\
         & s_y &  & -C_y s_y \\
         &  & s_z & -C_z s_z \\
         &  &  & 1
    \end{bmatrix}
\]

Dall'ultima riga della matrice notiamo che la coordinata $w$ viene sempre mantenuta a 1: in questo caso la perspective divide risulta quindi ininfluente.

\begin{figure}[ht]
	\centering
	\includegraphics[width=1\textwidth]{cap3/Proiezione_ortografica}
	\caption{Volume di vista della proiezione ortografica}
    \label{fig:proiezione_ortografica}
\end{figure}

\subsubsection{Proiezioni prospettiche}
Nelle proiezioni prospettiche il centro di proiezione si trova a distanza finita dal piano di proiezione: gli oggetti più vicini all'osservatore risulteranno più grandi e viceversa.
Il volume di vista è un tronco di piramide (visibile in \cref{fig:proiezione_prospettica}) che dovrà essere trasformato in parallelepipedo, in cui i punti del piano vicino fungono da riferimento e quelli del piano lontano vengono compressi in un rettangolo più piccolo: il parallelepipedo dovrà essere infine trasformato nel cubo normalizzato tramite una proiezione ortografica.

Detta $d$ la distanza del piano di proiezione dal centro di proiezione, vogliamo che i punti $P = \begin{bmatrix} x & y & z & 1 \end{bmatrix}$ vengano mappati in altri punti in cui la grandezza dipende dalla distanza: l'obiettivo è quindi costruire una matrice $M_{persp}$ tale che, dopo la perspective divide, si ottenga $x' = x \cdot d / z$ e $y' = y \cdot d /z$.

Per ottenere questa divisione per $\frac{z}{d}$, è sufficiente che la quarta riga della matrice produca $w = \frac{z}{d}$: la quarta riga della matrice è quindi $\begin{bmatrix} 0 & 0 & \frac{1}{d} & 0 \end{bmatrix}$.
Le prime due righe lasciano invece $x$ e $y$ inalterati, poiché la divisione per $w$ penserà a scalarle correttamente.

Rimane da determinare la terza riga della matrice, che deve rimappare $z$ nell'intervallo $[-1, 1]$ del cubo canonico.
Essa non può dipendere da $x$ e $y$ (altrimenti la profondità di un punto cambierebbe in base alla sua posizione orizzontale o verticale), assumendo quindi necessariamente la forma $\begin{bmatrix} 0 & 0 & a & b \end{bmatrix}$, dove $a$ e $b$ sono due incognite.
Notiamo che, dopo la perspective divide, la coordinata $z$ risultante è:
\[
    z_{NDC} = \frac{az + b}{z/d} = \frac{(az + b)d}{z}
\]
Imponendo le condizioni per cui il piano vicino deve essere mappato in $z = -1$ e il piano lontano in $z = +1$, e assumendo che il piano di proiezione corrisponda al piano vicino ($d = near$):
\[
    \begin{cases} 
        \dfrac{(a \cdot near + b) \cdot d}{near} = -1 \\[10pt]
        \dfrac{(a \cdot far + b) \cdot d}{far} = +1 
    \end{cases}
    \underset{d\ =\ near}{\implies}
    \quad
    \begin{cases} 
        a \cdot near + b = -1 \\[6pt]
        \dfrac{(a \cdot far + b) \cdot near}{far} = +1 
    \end{cases}
\]
La cui soluzione è:
\[
    a = \frac{far + near}{far - near}, \qquad b = -\frac{2 \cdot far \cdot near}{far - near}
\]
La matrice $M_{persp}$ è quindi:
\[
    M_{persp} = 
    \begin{bmatrix}
        1 &   &   &   \\
          & 1 &   &   \\
          &   & \frac{far + near}{far - near} & -\frac{2 \cdot far \cdot near}{far - near} \\
          &   & \frac{1}{near} & 0
    \end{bmatrix}
\]

Una volta trasformato il volume di vista in un parallelepipedo, è sufficiente applicare $P_{ortho}$ per normalizzare nel cubo NDC, ottenendo la matrice di proiezione prospettica completa:
\[
    P_{persp} = P_{ortho} \cdot M_{persp}
\]

\begin{figure}[H]
	\centering
	\includegraphics[width=1\textwidth]{cap3/Proiezione_prospettica}
	\caption{Volume di vista della proiezione prospettica}
    \label{fig:proiezione_prospettica}
\end{figure}

% ====================================================
% 3.3.4 - CLIPPING E VIEWPORT 
% ====================================================
\subsection{Clipping e viewport}
Una volta ottenuto il cubo in coordinate NDC, viene eseguito il \emph{clipping}, in cui le porzioni di primitive geometriche che non rientrano nel volume di vista vengono scartate dalla pipeline: questo permette di eseguire i calcoli successivi solo per i vertici che rientrano effettivamente nell'area da visualizzare.
Il risultato di quest'operazione è schematizzato in figura.

Infine i valori delle coordinate $x$ e $y$ passano dal sistema NDC al sistema di \emph{viewport}, ovvero in coordinate dell'effettiva area di visualizzazione in pixel, di cui è necessario conoscere l'origine del sistema di riferimento $(x_0, y_0)$ e l'altezza e larghezza dell'area $(height, width)$:
\[
    x_{vp} = \frac{width}{2} \cdot x_{NDC} + \left(x_0 + \frac{width}{2}\right)
\]
\[
    y_{vp} = \frac{height}{2} \cdot y_{NDC} + \left(y_0 + \frac{height}{2}\right)
\]

A questo punto della pipeline termina il sottosistema geometrico, per passare al sottosistema di rasterizzazione.

% ====================================================
% 3.3.5 - RASTERIZZAZIONE 
% ====================================================
\subsection{Rasterizzazione}
Il sottosistema raster è l'ultima fase della pipeline di rendering.
Il processo di rasterizzazione si occupa di convertire le primitive geometriche (punti, linee, poligoni) definite in coordinate di viewport nella loro rappresentazione discreta composta da frammenti: ne è riportato un esempio in \cref{fig:rasterizzazione}
Nel passaggio dai vertici ai frammenti, la GPU interpola tutti i dati dei vertici per ogni frammento grazie alle coordinate baricentriche, discusse in precedenza.
Dopo la produzione dei frammenti interviene il \emph{fragment shader}, utilizzato principalmente per calcolarne il colore.

\begin{figure}[ht]
	\centering
	\includegraphics[width=0.5\textwidth]{cap3/Rasterizzazione}
	\caption[Rasterizzazione di primitive geometriche]{Dettaglio di rasterizzazione di alcune primitive geometriche, in cui si mostra come alcuni pixel vengano ``accesi'' per approssimarne le forme}
    \label{fig:rasterizzazione}
\end{figure}

Una volta terminata l'elaborazione dei frammenti, tramite il \emph{depth test} la GPU decide se scartare il frammento -- ad esempio perché coperto da un oggetto più vicino alla camera -- o scriverlo nel color buffer. 
In quest'ultimo caso, il colore viene eventualmente miscelato a quello già presente qualora il \emph{blending} sia attivo.
Il blending è una tecnica utilizzata principalmente per implementare la trasparenza: invece di sovrascrivere il color buffer con il colore del frammento candidato, i due colori vengono miscelati in base al canale \emph{alpha} $\alpha \in [0, 1]$ del frammento in arrivo.
Un esempio di funzione di blending può essere la seguente, denominata da OpenGL come \texttt{GL\_ONE\_MINUS\_SRC\_ALPHA}:
\[
    C_{out} = \alpha \cdot C_{src} + (1 - \alpha) \cdot C_{dst}
\]

% ====================================================
%
% 3.4 - MODELLI DI SHADING E ILLUMINAZIONE
%
% ====================================================
\section{Modelli di illuminazione e shading}
I modelli di illuminazione descrivono i fattori che determinano il colore di una superficie in un determinato punto, tramite le interazioni tra luci e superfici, le proprietà delle superfici (materiali) e la natura della radiazione luminosa incidente.
Sono classificati in due gruppi, i modelli di illuminazione locale e i modelli di illuminazione globale.

Un modello di illuminazione è detto locale quando calcola l'intensità luminosa di un singolo punto della superficie considerando esclusivamente la luce diretta dalle sorgenti luminose.
Non vengono quindi calcolate ombre, riflessioni all'interno dell'ambiente e/o rifrazioni.

Un modello di illuminazione è invece detto globale quando calcola il colore in un punto considerando sia la luce diretta che quella che raggiunge il punto dopo essere stata riflessa e/o trasmessa dalle altre superfici presenti nell'ambiente circostante.
I modelli di illuminazione globale più conosciuti sono il \emph{ray tracing}, che simula il percorso dei raggi luminosi, e il \emph{radiosity}, che modella lo scambio di energia luminosa tra superfici tramite equazioni di conservazione dell'energia.

Sebbene i modelli locali offrano semplicità e velocità di calcolo, solo con i modelli globali si può ottenere il fotorealismo: le differenze tra i modelli sono visibili nelle immagini a confronto in \cref{fig:modelli_illuminazione}.

\begin{figure}[H]
    \centering
    \begin{subfigure}[b]{0.64\textwidth}
        \centering
        \includegraphics[width=\textwidth]{cap3/IlluminazioneLocale}
        \caption{Illuminazione locale}
    \end{subfigure}
    \hfill
    \begin{subfigure}[b]{0.33\textwidth}
        \centering
        \includegraphics[width=\textwidth]{cap3/IlluminazioneGlobale}
        \caption{Illuminazione globale}
    \end{subfigure}
    \caption[Differenza tra modelli di illuminazione locali e globali]{Immagini in cui si nota il divario tra i modelli di illuminazione locali e globali: è immediatamente evidente la differenza nella gestione delle ombre e delle riflessioni}
    \label{fig:modelli_illuminazione}
\end{figure}

Mentre il calcolo dell'illuminazione si occupa di definire l'intensità luminosa di un punto in base alle proprietà della luce e del materiale, lo \emph{shading} specifica come tale luce viene effettivamente utilizzata per colorare i pixel dell'immagine finale.

Nei paragrafi seguenti vengono analizzati il modello di Phong e la sua variante migliorata di Blinn-Phong, per poi illustrare come questi vengono applicati attraverso i principali modelli di shading.

% ====================================================
% 3.4.1 - MODELLO DI ILLUMINAZIONE DI PHONG
% ====================================================
\subsection{Modello di illuminazione di Phong}
Il modello di illuminazione di Phong approssima tutti i contributi della luce in 3 componenti -- ambientale, diffusivo e speculare -- e fa uso di sorgenti di luce puntiformi, ovvero sorgenti dotate di una posizione specifica nello spazio, prive di volume e area, che emettono luce uniformemente in ogni direzione.

\subsubsection{Componente ambientale}
Nella realtà gli oggetti che non sono direttamente colpiti dalla luce non appaiono completamente neri: questo perché vengono illuminati dalla luce riflessa dall'ambiente circostante.
Si approssimano tutti i contributi di questi riflessi nella luce ambientale, una luce diffusa e non direzionale che colpisce ugualmente tutte le superfici da tutte le direzioni.
La quantità di luce dell'ambiente riflessa dalla superficie dell'oggetto è determinata dal coefficiente di riflessione ambientale $k_a$: $k_a \in [0, 1]$ ed è una proprietà che caratterizza il materiale di cui la superficie è composta.

Sia $k_a$ il coefficiente di riflessione ambientale del materiale, e $I_a$ l'intensità della luce ambientale (supposta costante per tutti gli oggetti della scena), allora la componente ambientale $A$ è determinata da:
\[
    A = I_a \cdot k_a
\]
Gli oggetti illuminati da sola luce ambientale sono ancora uniformemente illuminati su tutta la superficie.

\subsubsection{Componente diffusiva}
Questo termine modella la riflessione lambertiana, tipica di materiali opachi e rugosi come il gesso.
Queste superfici appaiono ugualmente luminose da qualunque punto di vista vengano osservate poiché riflettono uniformemente la luce in tutte le direzioni.
La luminosità dipende solo dall'angolo $\theta$, ovvero l'angolo formato dalla direzione del raggio luminoso $L$ e la normale alla superficie nel punto d'incidenza $N$ (schematizzato in \cref{fig:vettori_diffusivi}): è quindi indipendente dalla posizione dell'osservatore.
La quantità di luce data dalla riflessione lambertiana della superficie dell'oggetto è determinata dal coefficiente di riflessione diffusiva $k_d \in [0, 1]$.

\begin{figure}[ht]
	\centering
	\includegraphics[width=0.4\textwidth]{cap3/VettoriDiffusiva}
	\caption{Vettori della componente diffusiva}
    \label{fig:vettori_diffusivi}
\end{figure}

Sia $k_d$ il coefficiente di riflessione diffusiva del materiale e $I_p$ l'intensità della sorgente luminosa puntiforme, allora la componente diffusiva $D$ è data da:
\[
    D = I_p \cdot k_d \cdot \cos{\theta}
\]
Nel caso in cui $N$ e $L$ siano normalizzati utilizziamo il prodotto scalare:
\[
    D = I_p \cdot k_d \cdot (\hat{N} \cdot \hat{L})
\]

\subsubsection{Componente speculare}
Nel caso in cui la superficie di un oggetto non sia completamente opaca, la luce non viene riflessa ugualmente in tutte le direzioni, ma si concentra in un intervallo ristretto.
La direzione di riflessione $R$ è geometricamente la direzione di incidenza della luce $L$ riflessa rispetto ad $N$.
$L, R, N$ sono vettori che giacciono sullo stesso piano, e per definizione l'angolo di riflessione $\theta_R$ è uguale all'angolo di incidenza della luce $\theta_L$.
Il vettore $V$ è la direzione di vista (ovvero il vettore dato dalla posizione della camera sottratto alla posizione del suo obiettivo), e l'angolo formato col vettore di riflessione è indicato come $\alpha$.
Tutte queste componenti sono rappresentate in \cref{fig:vettori_speculare}.

Nei materiali perfettamente speculari (come gli specchi), l'osservatore vede il riflesso solo se $\alpha = 0$, quindi solo se il vettore di vista è allineato con il vettore riflessione.
Phong trova un modo per simulare i materiali non perfettamente speculari: fa decadere rapidamente la luce riflessa all'aumentare di $\alpha$, elevando $\cos{\alpha}$ per un coefficiente $n$: si genera così un cono intorno alla direzione di riflessione (visibile in \cref{fig:cono_riflessione}), in cui il riflesso è comunque visibile.
Questo coefficiente $n$ è detto \emph{shininess} (lucentezza), o coefficiente di riflessione speculare, e controlla la dimensione e l'intensità delle macchie di luce che appaiono sulle superfici quando è colpita dalla luce.
Nella formula $\cos^n\alpha$, l'aumento di $n$ corrisponde a una riduzione della dimensione del cono di direzioni intorno al vettore $R$ in cui la luce riflessa rimane visibile: valori alti (100 - 500) corrispondono a un riflesso nitido e puntiforme, tipico delle superfici metalliche, mentre valori bassi ($<$ 100) producono un'ampia zona di massima lucentezza più sfocata, tipica ad esempio di plastica e cera.
La quantità di luce riflessa dalla superficie dell'oggetto è invece determinata dal coefficiente di riflessione speculare $k_s \in [0, 1]$.

\begin{figure}[ht]
    \centering
    \begin{subfigure}[b]{0.49\textwidth}
        \centering
        \includegraphics[width=\textwidth]{cap3/VettoriRiflessione}
        \caption{Vettore di riflessione}
    \end{subfigure}
    \hfill
    \begin{subfigure}[b]{0.49\textwidth}
        \centering
        \includegraphics[width=\textwidth]{cap3/ConoRiflessione}
        \caption{Vettore di vista, cono di riflessione}
        \label{fig:cono_riflessione}
    \end{subfigure}
    \caption[Vettori della componente speculare]{Vettori che concorrono alla creazione della componente speculare}
    \label{fig:vettori_speculare}
\end{figure}

Dati il coefficiente $k_s$, la shininess $n$, l'angolo $\alpha$ formato tra $R$ e $V$ e $I_p$ l'intensità della sorgente luminosa, allora la componente speculare $S$ è data da:
\[
    S = I_p \cdot k_s \cdot \cos^n\alpha
\]
Nel caso in cui $R$ e $V$ siano normalizzati:
\[
    S = I_p \cdot k_s \cdot (\hat{R} \cdot \hat{V})^n
\]

\subsubsection{Modello completo}
Il modello di illuminazione di Phong unisce le tre componenti appena studiate, ovvero ambientale $A$, diffusiva $D$ e speculare $S$, il cui contributo è mostrato in un esempio in \cref{fig:componenti_phong}.

\begin{figure}[ht]
	\centering
	\includegraphics[width=\textwidth]{cap3/ComponentiPhong}
	\caption[Componenti del modello di illuminazione di Phong]{Esempio di come ogni componente (ambientale, diffusiva, speculare) dia il suo contributo per la creazione dell'immagine finale}
    \label{fig:componenti_phong}
\end{figure}

Nel calcolo dell'illuminazione, si può voler tenere conto anche dell'attenuazione dell'intensità luminosa, ad esempio all'aumentare della sua distanza dal punto di incidenza $d_L$: $f_{att} = 1 / {d_L}^2$. 
Unendo tutti questi contributi si ha la formula $I = A + f_{att} (D + S)$, che espansa fornisce l'equazione del modello di illuminazione di Phong, in cui i vettori utilizzati sono schematizzati in \cref{fig:vettori_phong}:
\[
    I = I_a \cdot k_a + f_{att} \cdot I_p [k_d (\hat{N} \cdot \hat{L}) + k_s (\hat{R} \cdot \hat{V})^n]
\]
Nel caso siano presenti più luci è sufficiente sommare i termini per ogni sorgente puntiforme:
\[
    I = I_a \cdot k_a + \sum_{l \in luci} f_{att l} \cdot I_{pl} [k_d (\hat{N} \cdot \hat{L_l}) + k_s (\hat{R_l} \cdot \hat{V})^n]
\]

\begin{figure}[H]
	\centering
	\includegraphics[width=0.65\textwidth]{cap3/VettoriPhong}
	\caption{Vettori del modello di illuminazione di Phong}
    \label{fig:vettori_phong}
\end{figure}

% ====================================================
% 3.4.2 - MODELLO DI ILLUMINAZIONE DI BLINN-PHONG
% ====================================================
\subsection{Modello di illuminazione di Blinn-Phong}
Blinn propone una variante del modello di Phong per quanto riguarda la componente speculare.
Poiché il ricalcolo continuo del vettore riflessione $R$ è computazionalmente oneroso, introduce l'\emph{halfway vector}.

L'halfway vector $H$ è un vettore a metà strada tra la direzione di vista $V$ e la luce $L$: è quindi la diagonale del parallelogramma formato da $V$ ed $L$ (visibile in \cref{fig:vettori_blinnphong}).
Dato che sappiamo che la diagonale del parallelogramma formato da due vettori non è altro che la loro somma, il vettore $H$ si calcola:
\[
    H = L + V
\]
Spostando $V$ verso $R$ accade che $H$ si avvicina a $N$.
Nel punto di massima riflessione, ovvero $V = R$, si ha che $H$ è allineato ad $N$: l'angolo $\beta$, formato tra $N$ e $H$, approssima bene l'angolo $\alpha$.

Blinn decide quindi di sostituire $(\hat{R} \cdot \hat{V})^n$ con $(\hat{N} \cdot \hat{H})^n$, creando così l'equazione del modello di illuminazione di Blinn-Phong:
\[
    I = I_a \cdot k_a + f_{att} \cdot I_p [k_d(\hat{N} \cdot \hat{L}) + k_s (\hat{N} \cdot \hat{H})^n]
\]

Bisogna però notare che spesso $\beta$ è più piccolo di $\alpha$ usato in Phong, tranne quando la superficie è vista da un angolo molto ripido. Per avvicinarsi di più al risultato di $(R \cdot V)^n$ di Phong, in $(N \cdot H)^{n_\beta}$ sarebbe opportuno scegliere $n_\beta > n$.
Inoltre, mentre Phong produce sempre riflessi rotondi su una superficie piana, quelli di Blinn-Phong diventano ellittici quando la superficie viene vista da un angolo ripido, ad esempio quando il sole vicino all'orizzonte viene riflesso sulla superficie del mare.

\begin{figure}[ht]
	\centering
	\includegraphics[width=0.7\textwidth]{cap3/VettoriBlinnPhong}
	\caption{Costruzione dell'halfway vector}
    \label{fig:vettori_blinnphong}
\end{figure}

% ====================================================
% 3.4.3 - MODELLI DI SHADING
% ====================================================
\subsection{Modelli di shading}
I modelli di shading rappresentano il processo matematico e algoritmico utilizzato per determinare il colore di tutti i pixel che ricoprono una superficie, applicando un determinato modello di illuminazione.
Si differenziano principalmente per la frequenza con cui viene ricalcolato il modello di illuminazione, dando diversi effetti visibili in \cref{fig:modelli_shading}.

\begin{figure}[ht]
	\centering
	\includegraphics[width=1\textwidth]{cap3/Shading_models}
	\caption[Diversi modelli di shading a confronto]{Confronto tra i principali modelli di shading; lo shading interpolativo, non rappresentato, si colloca concettualmente tra flat e Gouraud shading}
    \label{fig:modelli_shading}
\end{figure}

\subsubsection{Flat shading}
È il modello più semplice ed efficiente: il modello di illuminazione è applicato una sola volta per poligono.
Questo può rendere le divisioni in facce visibili su superfici curve e crea una discontinuità netta tra facce adiacenti.
Viene ad esempio calcolato il colore solo per il centroide del triangolo, per poi applicarlo a tutto il poligono.

\subsubsection{Shading interpolativo}
L'equazione di illuminazione viene calcolata per ogni vertice del poligono, per poi eseguire un'interpolazione lineare per colorarne ogni pixel.
Poiché non teniamo conto che un triangolo fa parte di una geometria più complessa, abbiamo una sola normale per triangolo.
Questo causa brusche variazioni delle normali ai bordi dei poligoni adiacenti, rendendo i bordi nettamente visibili: questo fenomeno prende il nome di \emph{Mach banding}.

\subsubsection{Gouraud shading}
Rappresenta la diretta evoluzione dello shading interpolativo.
Esegue comunque il calcolo dell'illuminazione per ogni vertice del poligono, ma invece di utilizzare la normale del poligono usa quella della superficie che sta approssimando.
Se la normale della superficie non è nota, la si approssima con la media delle normali ai poligoni che condividono il vertice.
Si procede poi con l'interpolazione della luce calcolata nei vertici.
Non è ideale per materiali molto riflettenti o speculari, poiché i punti luce possono andare persi se ricadono all'interno di un poligono senza colpirne i vertici.

\subsubsection{Phong shading}
È il modello più sofisticato e realistico ma anche più computazionalmente costoso.
Prevede il calcolo dell'illuminazione per ogni pixel del poligono, utilizzando il vettore normale ottenuto dall'interpolazione delle normali note.
Promette così di gestire correttamente le riflessioni speculari, anche quando queste non colpiscono direttamente i vertici.

% ====================================================
% CONCLUSIONE CAPITOLO 3
% ====================================================
\bigskip
Disponendo ora degli strumenti necessari alla comprensione della pipeline grafica, è possibile procedere con la progettazione dell'architettura del software.